\section{QmisID: Charge Misidentification}

This paper is geared towards the goal of learning particle-level detector response functions that are fully differential. These encompass both continuous functions (i.e. resolution effects) as well as discrete functions. 
Here we focus on the second task: learning to predict whether a particle produced via pair-production has been misiedentified (e+ detected as e- or vice cerse) during reconstruction. 

This is very relevant to the two-lepton same-charge (2LSC) channel of the di-Higgs decay. This is one of the nine sub-channels in the di-Higgs multi-lepton decay, which consists of two light leptons (electrons or muons) of the same charge. QmisID events occur when when a reco label of same-charge (SC) is mistakenly assigned to a two-lepton event which was opposite-charge (OC) in truth. 
Two physical mechanisms can come into play. First, the curvature of a charged particle's trajectory in a uniform magnetic field asymptotically approaches zero at high transverse momentum. Minor deflection of the track, for instance due to scattering inside the detector, can lead to a charge flip. 

The second mechanism is bremsstrahlung followed by photon conversion:
an electron interacts with the detector material, radiates a photon,
and that photon may itself convert into an \(e^{+}e^{-}\) pair.
Above \(\sim\!10\)~MeV, bremsstrahlung is the dominant energy-loss
mechanism \cite{wichmann2008diplomarbeit}. Charge misassignment then occurs when one conversion
electron carries nearly all of the original photon energy and the
reconstruction algorithm follows that secondary track instead of the
original one; the probability of this grows with photon
energy~\cite{wichmann2008diplomarbeit}. The effect has a strong functional dependence on pseudorapidity \(\eta\), as the density of the inner detector increases from 0 to 2.5 in absolute value of \(\eta\). Pseudorapidity \(\eta\) is a spatial coordinate, where \(\theta\) is the polar angle between the particle's momentum vector and the beam axis.

\begin{equation}
    \eta = -\ln\!\left(\tan\frac{\theta}{2}\right).
\end{equation}

We would like to learn to predict the probability of charged particle, including electron, misidentification with a method that is fully differential.